\chapter{Introduction}
\label{ch:intro}

Electronic documents with sensitive information are used across nearly every industry and profession. Examples include health records, legal counsel, financial data, trade secrets, and employment contracts. These sensitive documents are often stored and shared using cloud services like managed file transfer (MFT) applications because while security and data protection are of critical importance, the files must also be easily accessible and sharable in order to be useful. Users trust these platforms to keep their data secure. Documents remain an important means of digital transmission deserving robust cybersecurity protections.

However, in recent years, a series of hacking incidents targeting cloud-based file transfer applications compromised millions of confidential documents, including personal health data, government records, and legal files~\cite{kondruss_moveit_2023,cisa_cl0p_2023}. The largest incident in 2023 resulted in the exposure of 95{,}788{,}491 individuals' personal information~\cite{simas_unpacking_2023}. Worse, these incidents represent an ongoing pattern. Despite mass exploit events, the same attack pattern has continued unabated as it exposes a systematic design flaw with how many documents are protected today. From 2020 to 2025, more than eight vulnerabilities allowing remote privileged access to MFTs were discovered and abused, resulting in major data breaches (see Section~\ref{sec:mft-hacks} for further discussion)~\cite{cisa_cl0p_2023}. Although MFT applications encrypt files ``at rest,'' all documents are unencrypted and readable by a privileged user session. Once an attacker gains access every document on the server can be exfiltrated. The current security model fails because it relies on perimeter defense rather than per-document protection. There is a pressing need to understand the limitations of existing document protection systems and develop novel systems to provide better security.

While much research and development has been performed to provide confidentiality for web traffic (HTTPS, TLS, etc.) and instant messaging (end-to-end encrypted Signal, WhatsApp, etc.), document-level encryption has received comparatively little attention. From a usability perspective, the state of document encryption is troublesome. Password-based encryption built into tools like Microsoft Word and Adobe Acrobat today works nearly the same as it did when introduced in the 1990s~\cite{adobe_pdf_timeline}: set a password and hope you don't forget it. It has major usability issues such as the inability to recover or reset a forgotten password. Password-based file encryption is also vulnerable to brute-force attacks~\cite{danczul_cuteforce_2013,hranicky_experimental_2024} and unauthorized sharing. The result is that despite the availability of encryption features, many sensitive documents remain inadequately protected in practice. A successful document encryption system must provide adequate usability, security, and deployability; without all three, adoption will not follow.

Significant advances in passwordless authentication---notably FIDO2 and WebAuthn, developed by the FIDO Alliance and W3C---provide compelling components for a novel document encryption method that offers strong, modern protection without compromising usability. As of 2026, FIDO2 and WebAuthn have begun to cross the deployability and adoption barriers that have historically prevented password replacement mechanisms from gaining traction~\cite{fido_state_of_passkeys_2026}. To date, FIDO2 and WebAuthn applications have focused primarily on user authentication (the average consumer's user login to a website). However, the potential benefits extend beyond user authentication. The WebAuthn PRF extension enables end-to-end encryption derived from a passkey. While the WebAuthn PRF extension has seen early use in password managers and cryptocurrency wallets, no prior work has applied FIDO2 and WebAuthn to perform document encryption where files must be shared between parties, nor has any academic study evaluated the feasibility, deployment, and usability of such an approach.

This research bridges that gap. We first conducted a systematic analysis of existing document protection methods using a structured comparative framework adapted from usable security research methods for evaluating password alternatives~\cite{bonneau_quest_2012} and secure email~\cite{borisov_sok_2021}. We applied this framework to compare nine document protection methods across 15 design properties, identifying the benefits and limitations of current approaches. From this analysis, we derived design requirements for a system that is not password-based, not proprietary, does not require recipients to install software, and encrypts documents individually. We then designed and implemented a novel document encryption system leveraging FIDO2 passkeys, the WebAuthn PRF extension, and open-source age encryption to provide end-to-end protection of digital documents without passwords. Finally, we evaluated this system through a within-subjects user study (N=21) comparing the passwordless approach to traditional password-based encryption as provided by Adobe Acrobat, measuring usability and user acceptance with established instruments alongside qualitative interview data. Results show that the passwordless prototype achieved usability parity with Adobe Acrobat on first use (SUS 72.0 vs.\ 73.3, $p=.944$), while being rated significantly more useful by participants (VdL $p=.014$). When asked which method they would prefer for sharing multiple documents, 19 of 21 participants chose the passwordless approach, citing the password sharing burden and trust in biometric authentication as key factors.

This study is guided by the following research questions:

\begin{description}
    \item[R1:] What usability, deployability, and security limitations exist in current document encryption methods?
    \item[R2:] How can open-source protocols be used to design a document protection system that is passwordless, non-proprietary, and provides per-file encryption?\footnote{(1) \emph{Passwordless} meaning: no password or memory aspect needed for the scheme to work. (2) \emph{Non-proprietary} meaning: free, open-source and when files are shared, the receiver doesn't need custom software to view or interact with the file. (3) \emph{Per-file encryption} meaning: individual document files are encrypted uniquely and not all accessible at once.}
    \item[R3:] What is the user acceptance and usability of a passwordless document encryption system?
\end{description}

The main contributions of this work are as follows. First, to our knowledge this is the first study to apply FIDO2 and WebAuthn PRF to end-to-end document encryption rather than user authentication, and one of the first human-subjects usability evaluations of the PRF extension in any context. Second, we provide analysis of the pattern of MFT hacks and limitations of document encryption methods which grounds the research in real-world deployability, addressing a pressing problem demonstrated by millions of individuals affected by MFT data breaches. Third, we present a structured comparative framework for evaluating document encryption methods. Fourth, we provide empirical evidence on the usability and user acceptance of passwordless document encryption through a controlled lab study.

\section{Overview of Thesis}
\label{sec:thesis-overview}

Chapter~\ref{ch:intro} gave an introduction and explained the problem and research questions motivating this research. Chapter~\ref{ch:background} provides a summary of recent MFT hacking incidents and related work in existing literature to give context and background to the research. Chapter~\ref{ch:methodology} outlines the methodology of the research which is divided into three phases. Chapter~\ref{ch:evaluation} details the method and results of careful analysis and comparison of existing document encryption methods. Chapter~\ref{ch:system-design} explains the design requirements and process of developing the passwordless document encryption system. Chapter~\ref{ch:user-study} reviews the user testing performed including the IRB approval, screening criteria, tasks performed, and results gathered through both quantitative scales and qualitative interview questions. Chapter~\ref{ch:results} presents the results of the user study. Chapter~\ref{ch:discussion} presents discussion, limitations, and concluding perspectives.
